\begin{longtable}{
    %|>{\raggedright\arraybackslash}p{0.27\textwidth}
    |>{\raggedright\arraybackslash}p{0.23\textwidth}
    |p{0.55\textwidth}
    |p{0.32\textwidth}
    |p{0.27\textwidth}|
    } % 5 - 1.35 \ 
\caption{Классификация терминов цикла по общему содержанию}
\label{tab:cycle-classification-general} \\
\LongTableHeader{
    \hline
    \rowcolor{black}
    %\textcolor{white}{\textbf{Русский термин}} &
    \textcolor{white}{\textbf{Термин}} &
    \textcolor{white}{\textbf{Описание}} &
    \textcolor{white}{\textbf{Источник}} &
    \textcolor{white}{\textbf{Цитата из источника}} \\
    \hline
}

\longcaptioneachpage{Источник: составлено автором}

% Деловой цикл / Бизнес-цикл
\begin{minipage}[t]{\linewidth}
    \raggedright
    \begin{itemize}
        \item \textbf{деловой цикл / бизнес-цикл (рус.)}
        \item business cycle (англ.)
    \end{itemize}
\end{minipage}
&
Квалификатор \textquotedbl{}business\textquotedbl{} означает \textquotedbl{}дело\textquotedbl{}, \textquotedbl{}предпринимательство\textquotedbl{}, \textquotedbl{}хозяйственная деятельность\textquotedbl{} в широком институциональном смысле. Митчелл сознательно предпочёл это слово слову \textquotedbl{}trade\textquotedbl{}: он хотел подчеркнуть, что цикличность не сводится к торговым колебаниям, но укоренена в самой природе капиталистической бизнес-организации. Термин впервые зафиксирован у Джонсона (1887) в контексте оценки хозяйственных результатов за период достаточной длины. Бёрнс и Митчелл в 1946 г. дали классическое определение: совокупные флуктуации экономической активности, проходящие через последовательные фазы расширения, пика, сжатия и дна. Именно Митчелл стандартизировал термин в книге \textquotedbl{}Business Cycles\textquotedbl{} (1913, NBER) и закрепил его как основу американской традиции. Besomi иронично замечает, что этот термин несёт в себе институционалистское убеждение Митчелла о том, что ритмические колебания неотделимы от \textquotedbl{}нашей нынешней бизнес-организации\textquotedbl{}. Сегодня \textquotedbl{}business cycle\textquotedbl{} – наиболее распространённый термин в мировой экономической литературе.
&
\begin{minipage}[t]{\linewidth}
    \vspace{-7pt}
    \RaggedRight
    \begin{enumerate}
        \item Johnson J. (1887). \textit{North American Review}\cite{Johnson-1887} (первое употребление)
        \item Mitchell W. C. (1927). \textit{Business Cycles: The Problem and Its Setting}\cite{Mitchell-1927} (стандартизация)
    \end{enumerate}
\end{minipage}
&
\textquotedbl{}In order to evaluate the results of a business a period long enough to include a business cycle must be covered\textquotedbl{}
\begin{flushright}
    Johnson J., 1887\cite{Johnson-1887}
\end{flushright}
\\ \hline

% Экономический цикл
\begin{minipage}[t]{\linewidth}
    \raggedright
    \begin{itemize}
        \item \textbf{экономический цикл (рус.)}
        \item economic cycle (англ.)
        \item cycle économique (фр.)
        \item ciclo económico (исп.)
        \item ciclo economico (ит.)
    \end{itemize}
\end{minipage}
&
Квалификатор \textquotedbl{}economic\textquotedbl{} / \textquotedbl{}économique\textquotedbl{} указывает на широту охвата – всю совокупность экономических отношений, а не только бизнес-сферу или торговлю. В английском языке термин менее распространён, чем \textquotedbl{}business cycle\textquotedbl{}, однако в заглавиях крупных работ появился уже в 1914 г. (Moore) и 1917 г. (J. M. Clark). Термин доминирует во французской (cycle économique), итальянской (ciclo economico) и испанской (ciclo económico) научных традициях, тогда как в советской традиции он был фактически основным наряду с \textquotedbl{}хозяйственным циклом\textquotedbl{}. По сравнению с \textquotedbl{}деловым циклом\textquotedbl{} он несёт более широкий охват: включает как рыночные, так и нерыночные секторы хозяйства, что делало его удобным для советских авторов, описывавших явления в капиталистических странах. Romer (1993) критиковала само понятие \textquotedbl{}цикл\textquotedbl{} как подразумевающее ритмическую регулярность, которой флуктуации на деле не обнаруживают. В современной литературе \textquotedbl{}economic cycle\textquotedbl{} используется как более нейтральное, широкое понятие, взаимозаменяемое с \textquotedbl{}business cycle\textquotedbl{} в международном контексте.
&
\begin{minipage}[t]{\linewidth}
    \vspace{-7pt}
    \RaggedRight
    \begin{enumerate}
        \item Moore H. L. (1923). \textit{Generating Economic Cycles}\cite{Moore-1923}
    \end{enumerate}
\end{minipage}
&
\textquotedbl{}These features in the theory and history of tides will probably be found to have their counterparts in the theory and history of economic cycle\textquotedbl{}
\begin{flushright}
    Moore H. L., 1923\cite{Moore-1923}
\end{flushright}
\\ \hline

% Конъюнктурный цикл
\begin{minipage}[t]{\linewidth}
    \raggedright
    \begin{itemize}
        \item \textbf{конъюнктурный цикл (рус.)}
        \item konjunkturzyklus (нем.)
        \item conjoncture (фр.)
        \item congiuntura (ит.)
    \end{itemize}
\end{minipage}
&
\textquotedbl{}Konjunktur\textquotedbl{} происходит от латинского conjungere (соединять) и первоначально обозначал совокупность рыночных условий в данный момент, без какой-либо импликации цикличности – астрологическое по происхождению слово, обозначавшее \textquotedbl{}стечение обстоятельств\textquotedbl{}. Термин был введён Лассалем и принят Шеффле и Вагнером с целью подчеркнуть, что совокупность рыночных условий существует независимо от воли и устремлений отдельных хозяйственных агентов. Лишь позднее \textquotedbl{}konjunktur\textquotedbl{} поглотил идею ритмического изменения, а производный составной термин \textquotedbl{}konjunkturzyklus\textquotedbl{} появился в немецкоязычной литературе в 1933 г. Немецкоязычная традиция различает \textquotedbl{}konjunktur\textquotedbl{} (текущее состояние хозяйственной конъюнктуры) и \textquotedbl{}konjunkturzyklus\textquotedbl{} (полный цикл её изменения), тогда как в русском переводе оба понятия объединяются под словом \textquotedbl{}конъюнктурный цикл\textquotedbl{}. Феллден (1927) характеризовал этот цикл как \textquotedbl{}чередующееся изменение между процветанием и депрессией\textquotedbl{}. Bank-Lexikon (1988) уточняет, что \textquotedbl{}konjunktur\textquotedbl{} относится не к конкретной ситуации (росту или кризису), а к циклическому процессу в целом.
&
\begin{minipage}[t]{\linewidth}
    \vspace{-7pt}
    \RaggedRight
    \begin{enumerate}
        \item Palyi M., Quittner P. (1933). \textit{Handwörterbuch des Bankwesens}\cite{Palyi-1933} (первое употребление)
        \item Hohlstein M. et al. (1994). \textit{Lexikon der Volkswirtschaft}\cite{Hohlstein-1994}
    \end{enumerate}
\end{minipage}
&
\textquotedbl{}The whole oscillation process of economic development through certain economic phases is called konjunkturzyklus\textquotedbl{}
\begin{flushright}
    Hohlstein M. et al., 1994\cite{Hohlstein-1994}
\end{flushright}
\vspace{7pt}
\textquotedbl{}The term konjunktur as introduced by Lassalle and adopted by Schäffle and Wagner was meant to emphasize the independence of the totality of market conditions from the actions and aspirations of the economic agents\textquotedbl{}
\begin{flushright}
    Kuznets S., 1930\cite{Kuznets-1930}
\end{flushright}
\\ \hline

% Коммерческий цикл
\begin{minipage}[t]{\linewidth}
    \raggedright
    \begin{itemize}
        \item \textbf{коммерческий цикл (рус.)}
        \item commercial cycle (анг.)
    \end{itemize}
\end{minipage}
&
Квалификатор \textquotedbl{}commercial\textquotedbl{} – \textquotedbl{}торговый\textquotedbl{}, \textquotedbl{}коммерческий\textquotedbl{} – восходит к эпохе, когда в центре внимания экономистов находились именно торговые и кредитные отношения, а не промышленный капитал. Термин появился в 1833 г. у Джона Вейда, английского журналиста, и представлял собой первое документально подтверждённое обозначение регулярно повторяющегося хозяйственного цикла. Митчелл (1927) и Besomi зафиксировали это первенство. Уже Фуллартон (1844) использовал \textquotedbl{}commercial cycle\textquotedbl{} в банковском контексте, а анонимные авторы 1848 и 1866 гг. – для описания регулярных волн \textquotedbl{}возбуждения\textquotedbl{} и \textquotedbl{}реакции\textquotedbl{} в коммерческом мире. К 1920-м гг. термин вышел из употребления, уступив место \textquotedbl{}trade cycle\textquotedbl{} в британской традиции и \textquotedbl{}business cycle\textquotedbl{} в американской. Его историческое значение – в том, что он зафиксировал первую попытку концептуализировать цикличность как закономерность, а не как патологию.
&
\begin{minipage}[t]{\linewidth}
    \vspace{-7pt}
    \RaggedRight
    \begin{enumerate}
        \item Wade J. (1833). \textit{History of the Middle and Working Classes}\cite{Wade-1833} (первое употребление)
    \end{enumerate}
\end{minipage}
&
\textquotedbl{}The commercial cycle is ordinarily completed in five or seven years, within which terms it will be found … alternate periods of prosperity and depression been experienced\textquotedbl{}
\begin{flushright}
    Wade J., 1833\cite{Wade-1833}
\end{flushright}
\\ \hline

% Хозяйственный цикл
\begin{minipage}[t]{\linewidth}
    \raggedright
    \begin{itemize}
        \item \textbf{хозяйственный цикл (рус.)}
        \item wirtschaftszyklus (нем.)
    \end{itemize}
\end{minipage}
&
Квалификатор \textquotedbl{}хозяйственный\textquotedbl{} – нем. Wirtschaft – обозначает более широкое понятие, чем \textquotedbl{}деловой\textquotedbl{} или \textquotedbl{}коммерческий\textquotedbl{}: он охватывает всю совокупность производственно-распределительной деятельности, включая натуральное хозяйство, государственный сектор и плановую экономику. Немецкий эквивалент \textquotedbl{}wirtschaftszyklus\textquotedbl{} тесно связан с понятием \textquotedbl{}wirtschaftliche konjunkturen\textquotedbl{}, разрабатывавшимся в немецкоязычной политической экономии XX в. В советской и постсоветской традиции \textquotedbl{}хозяйственный цикл\textquotedbl{} использовался наряду с \textquotedbl{}экономическим циклом\textquotedbl{} как его синоним с несколько более широким охватом. Понятие исторически связано с \textquotedbl{}промышленным циклом\textquotedbl{} (industriekreislauf) Маркса (1867). Besters (1963) определяет конъюнктурные колебания в рамках этой традиции через отклонение национального дохода от уровня, необходимого для поддержания полной занятости. Термин подчёркивает воспроизводственный характер цикла – цикл как возобновление хозяйственного кругооборота.
&
\begin{minipage}[t]{\linewidth}
    \vspace{-7pt}
    \RaggedRight
    \begin{enumerate}
        \item Besters H. (1963). \textit{Wirtschaftliche Konjunkturen}\cite{Besters-1963}
    \end{enumerate}
\end{minipage}
&
\textquotedbl{}The term 'wirtschaftszyklus' indicates fluctuations of national income around the level required to maintain the full employment of the national productive capacity\textquotedbl{}
\begin{flushright}
    Besters H., 1963\cite{Besters-1963}
\end{flushright}
\\ \hline

% Торговый цикл
\begin{minipage}[t]{\linewidth}
    \raggedright
    \begin{itemize}
        \item \textbf{торговый цикл (рус.)}
        \item trade cycle (анг.)
    \end{itemize}
\end{minipage}
&
Квалификатор \textquotedbl{}trade\textquotedbl{} – \textquotedbl{}торговля\textquotedbl{}, \textquotedbl{}торговые отношения\textquotedbl{} – отражает британскую традицию, в которой цикличность ассоциировалась прежде всего с движением торговли, кредита и банковских резервов, а не с промышленной организацией в целом. Первое употребление \textquotedbl{}trade cycle\textquotedbl{} как самостоятельного сочетания зафиксировано в 1857 г. (Anonymous), хотя уже в 1849 г. Харви использовал \textquotedbl{}cycle of trade\textquotedbl{}. The Economist употреблял это выражение в 1877 г. как устоявшееся, а Journal of the Statistical Society of London – в 1879 г. Окончательно \textquotedbl{}trade cycle\textquotedbl{} стал стандартным термином британской политической экономии после монографии Ф. Лавингтона (1922), которая, по словам Besomi, \textquotedbl{}установила это выражение в качестве стандартной британской терминологии\textquotedbl{}. В отличие от \textquotedbl{}business cycle\textquotedbl{}, термин не несёт институционалистской нагрузки Митчелла; он более узок по коннотациям. После Второй мировой войны \textquotedbl{}trade cycle\textquotedbl{} постепенно вытеснялся «business cycle» даже в британской литературе.
&
\begin{minipage}[t]{\linewidth}
    \vspace{-7pt}
    \RaggedRight
    \begin{enumerate}
        \item Lavington F. (1922). \textit{The Trade Cycle: an Account of the Causes Producing Rhythmical Changes in the Activity of Business}\cite{Lavington-1922} (стандартизация)
    \end{enumerate}
\end{minipage}
&
\textquotedbl{}This brief account of the trade cycle is intended to be a statement in as simple a form as possible of the main causes which underlie the rhythmical variations in the activity of business\textquotedbl{}
\begin{flushright}
    Lavington F., 1922\cite{Lavington-1922}
\end{flushright}
\\ \hline

% Промышленный цикл
\begin{minipage}[t]{\linewidth}
    \raggedright
    \begin{itemize}
        \item \textbf{промышленный цикл (рус.)}
        \item industrial cycle (анг.)
        \item industriezyklus (нем.)
    \end{itemize}
\end{minipage}
&
Квалификатор \textquotedbl{}industrial / industriell\textquotedbl{} специфицирует источник цикличности: речь идёт о цикле именно производственного, промышленного капитала. Маркс в \textquotedbl{}Капитале\textquotedbl{} (1867) описал последовательность «стагнация → умеренная активность → процветание → перепроизводство → кризис» как \textquotedbl{}цикл промышленных периодов\textquotedbl{} (the cycle of industrial periods) и связал его с периодическим обновлением основного капитала. Второй том \textquotedbl{}Капитала\textquotedbl{} (ред. Энгельса, 1885) развил понятие industriekreislauf как кругооборота промышленного капитала. В советской и постсоветской традиции \textquotedbl{}промышленный цикл\textquotedbl{} нередко использовался как синоним \textquotedbl{}экономического цикла\textquotedbl{}, хотя строго говоря относится к производственному сектору. В западной литературе термин встречается реже, уступив место \textquotedbl{}business cycle\textquotedbl{}, однако в историко-экономических работах о XIX в. остаётся незаменимым для обозначения марксовой концепции.
&
\begin{minipage}[t]{\linewidth}
    \vspace{-7pt}
    \RaggedRight
    \begin{enumerate}
        \item Marx K. (1867). \textit{Das Kapital. Vol. I}\cite{Marx-1867}
        \item Marx K. (1885). \textit{Das Kapital. Vol. II}\cite{Marx-1885}
    \end{enumerate}
\end{minipage}
&
\textquotedbl{}The cycle of industrial periods\textquotedbl{}
\begin{flushright}
    Marx K., 1867\cite{Marx-1867}
\end{flushright}
\vspace{7pt}
\textquotedbl{}Stagnation, moderate activity, prosperity, overproduction, crisis\textquotedbl{}
\begin{flushright}
    Marx K., 1885\cite{Marx-1885}
\end{flushright}
\\ \hline

\end{longtable}